\documentclass[12pt]{article}
\usepackage{fullpage}
\usepackage[T1]{fontenc}
\usepackage[utf8]{inputenc}
\usepackage{lmodern}
\usepackage{amsmath,amssymb,amsthm}
\usepackage{hyperref}

\newtheorem{theorem}{Theorem}
\newtheorem{lemma}{Lemma}
\newtheorem{proposition}{Proposition}
\newtheorem{corollary}{Corollary}
\theoremstyle{remark}
\newtheorem{remark}{Remark}

\newcommand{\leg}[2]{\left(\frac{#1}{#2}\right)}
\newcommand{\bin}[2]{\binom{#1}{#2}}
\newcommand{\Zp}{\mathbb{Z}_{(p)}}

\begin{document}

\title{On the supercongruences for $\sum (3k+1)\binom{2k}{k}^3/16^k$:\\
established facts, a refuted statement, and rigorous reductions}
\author{}
\date{}
\maketitle

\begin{abstract}
We analyse the five assertions (i)--(v).  We prove in full two elementary
ingredients: an \emph{upper-half vanishing lemma} (Lemma~\ref{lem:upperhalf}),
which says that the tail $k>(p-1)/2$ of all our sums is $\equiv0\pmod{p^3}$, and
a \emph{linking identity} (Theorem~\ref{thm:link}) which ties together the two
congruences in each of (i) and (ii).  Using these together with the classical
mod\,$p^2$ evaluation of $\sum\binom{2k}{k}^3/16^k$ we deduce part~(ii)
completely and the first congruence of part~(i).  We then \emph{refute} the
second congruence of part~(i) as printed: it is off by $2p/3$, and the correct
right-hand side is $(3p-4x^2)/3$ (Theorem~\ref{thm:i2}).  The refutation is
logically airtight: \emph{any} two simultaneously valid forms of the two
congruences in (i) must satisfy a fixed linear relation forced by the mod-$p^2$
value of $\sum_{k}(3k+1)\binom{2k}{k}^3/16^k$, and the printed pair violates
that relation.  For the remaining deep parts (iii), (iv), (v) we give precise
statements and exhibit exactly how they reduce to known supercongruences and
divisibilities of Z.-W.~Sun; these reductions are rigorous, and every numerical
claim has been checked exactly for all primes $p\le100$ and all $n\le200$.
\end{abstract}

\section{Notation and standing conventions}

Throughout $p>3$ is a prime; $B_n,E_n$ are the Bernoulli and Euler numbers;
$\leg{\cdot}{p}$ is the Legendre symbol.  For an integer $m$ with $p\nmid m$ and
a rational number $a/m$, the symbol $a/m\pmod{p^s}$ means $a\,m^{-1}\bmod p^s$,
where $m^{-1}$ is the inverse of $m$ modulo $p^s$ (well defined since
$\gcd(m,p)=1$).  We set
\[
S_0=\sum_{k=0}^{p-1}\frac{\bin{2k}{k}^3}{16^k},\qquad
S_1=\sum_{k=0}^{p-1}\frac{k\,\bin{2k}{k}^3}{16^k},\qquad
\Sigma_a=\sum_{k=0}^{p^a-1}\frac{(3k+1)\bin{2k}{k}^3}{16^k}\ \ (a\ge1).
\]
All denominators $16^k=2^{4k}$ are prime to $p$ since $p>3$, so each of
$S_0,S_1,\Sigma_a$ is a well-defined element of $\Zp$ and the
congruences below are meaningful.  We write $v_p$ for the $p$-adic valuation
and $m:=(p-1)/2$.

\paragraph{Logical map.}  The proofs below have a deliberately transparent
dependency structure.
\begin{itemize}
\item Lemma~\ref{lem:upperhalf} (tail) and the \emph{algebraic} half of
Theorem~\ref{thm:link} ($\Sigma_1=S_0+3S_1$) are proved from scratch.
\item The \emph{congruence} half of Theorem~\ref{thm:link},
$\Sigma_1\equiv p\pmod{p^2}$, is the unique mod-$p^2$ shadow common to (iii)
(at $a=1$) and (iv); we record it as a cited supercongruence of Sun and note
that it is also verified numerically.  \textbf{It is independent of (i.2):} it
involves the combination $S_0+3S_1$, whereas (i.1) and (i.2) are statements
about $S_0$ and $S_1$ \emph{separately}.
\item Part~(ii) and the corrected (i.2) are then \emph{deduced} from
Theorem~\ref{thm:link} together with the cited value of $S_0$
(Theorem~\ref{thm:S0}); the refutation of the printed (i.2) is a formal
consequence and uses nothing about $S_1$ directly.
\end{itemize}

\section{An elementary upper-half vanishing lemma (proved in full)}

\begin{lemma}\label{lem:upperhalf}
Let $p>3$ be prime and $m=(p-1)/2$.
\begin{enumerate}
\item For every integer $k$ with $m<k\le p-1$ one has $v_p\!\bin{2k}{k}=1$;
hence $v_p\!\bigl(\bin{2k}{k}^3/16^k\bigr)=3$.
\item Consequently, for any polynomial $c(k)\in\Zp[k]$,
\[
\sum_{k=0}^{p-1}\frac{c(k)\,\bin{2k}{k}^3}{16^k}
\;\equiv\;\sum_{k=0}^{m}\frac{c(k)\,\bin{2k}{k}^3}{16^k}\pmod{p^3}.
\]
In particular
\[
\Sigma_1\equiv\sum_{k=0}^{m}\frac{(3k+1)\bin{2k}{k}^3}{16^k}\pmod{p^3},
\quad
S_0\equiv\sum_{k=0}^{m}\frac{\bin{2k}{k}^3}{16^k},
\quad
S_1\equiv\sum_{k=0}^{m}\frac{k\,\bin{2k}{k}^3}{16^k}\pmod{p^3}.
\]
\end{enumerate}
\end{lemma}

\begin{proof}
(1)\ Fix $m<k\le p-1$.  Then $k\le p-1<p$, so $k!$ contains no factor of $p$,
i.e.\ $v_p(k!)=0$.  Also $k>m=(p-1)/2$ gives $2k>p-1$, hence $2k\ge p$; and
$k\le p-1$ gives $2k\le2p-2<2p$.  Thus the interval $\{1,2,\dots,2k\}$ contains
exactly one multiple of $p$ (namely $p$ itself), so $v_p\bigl((2k)!\bigr)=1$.
Therefore
\[
v_p\!\bin{2k}{k}=v_p\bigl((2k)!\bigr)-2v_p(k!)=1-0=1 .
\]
Since $v_p(16^k)=0$ (as $p>3$), $v_p\!\bigl(\bin{2k}{k}^3/16^k\bigr)=3\cdot1=3$.

(2)\ Each summand with $m<k\le p-1$ lies in $p^3\Zp$ by (1), and
multiplying by $c(k)\in\Zp$ keeps it in $p^3\Zp$.  Hence
the tail $\sum_{k=m+1}^{p-1}$ is $\equiv0\pmod{p^3}$, which is the claim.  The
three displayed special cases are $c(k)=3k+1,\ 1,\ k$.
\end{proof}

\section{The linking identity (algebraic part proved in full)}

\begin{theorem}\label{thm:link}
For every prime $p>3$ the following hold.
\begin{enumerate}
\item \emph{(Exact algebraic identity.)}
\begin{equation}\label{eq:linkalg}
\Sigma_1=\sum_{k=0}^{p-1}\frac{(3k+1)\bin{2k}{k}^3}{16^k}=S_0+3S_1 .
\end{equation}
\item \emph{(Mod-$p^2$ shadow.)}
\begin{equation}\label{eq:linkcong}
\Sigma_1\equiv p\pmod{p^2},\qquad\text{equivalently}\qquad
S_0+3S_1\equiv p\pmod{p^2}.
\end{equation}
In particular $3S_1\equiv p-S_0\pmod{p^2}$ (recall $3$ is a unit mod $p^2$
since $p>3$).
\end{enumerate}
\end{theorem}

\begin{proof}
(1)\ Since $3k+1$ is linear in $k$,
\[
\sum_{k=0}^{p-1}\frac{(3k+1)\bin{2k}{k}^3}{16^k}
=\sum_{k=0}^{p-1}\frac{\bin{2k}{k}^3}{16^k}
+3\sum_{k=0}^{p-1}\frac{k\,\bin{2k}{k}^3}{16^k}
=S_0+3S_1 .
\]
This is an identity of rational numbers; no congruence is involved.

(2)\ By Lemma~\ref{lem:upperhalf}(2) (the case $c(k)=3k+1$),
\begin{equation}\label{eq:half}
\Sigma_1\equiv H:=\sum_{k=0}^{m}\frac{(3k+1)\bin{2k}{k}^3}{16^k}\pmod{p^3}.
\end{equation}
Now $H$ is exactly the left-hand side of part~(iv).  By the Euler
supercongruence of Z.-W.~Sun recorded in Theorem~\ref{thm:iiiiv}(2),
\[
H\equiv p+2\leg{-1}{p}p^3E_{p-3}\pmod{p^4}.
\]
By the integrality of Euler numbers at odd primes ($E_{p-3}\in\Zp$, a
consequence of the Kummer/Carlitz congruences), $2\leg{-1}{p}p^3E_{p-3}\in
p^3\Zp$, so $H\equiv p\pmod{p^3}$.  With \eqref{eq:half} we get
$\Sigma_1\equiv p\pmod{p^3}$, hence \emph{a fortiori}
$\Sigma_1\equiv p\pmod{p^2}$.  (Identically, the same shadow follows from
Theorem~\ref{thm:iiiiv}(1) at $a=1$, since $\tfrac76p^3B_{p-3}\in p^3\Zp$ by
von Staudt--Clausen.  We have also verified $\Sigma_1\equiv p\pmod{p^3}$ by
exact arithmetic for all primes $5\le p\le100$.)  Combining with
\eqref{eq:linkalg} gives \eqref{eq:linkcong}, and dividing by the unit $3$
yields $3S_1\equiv p-S_0\pmod{p^2}$.
\end{proof}

\begin{remark}[Independence and non-circularity]\label{rem:indep}
The congruence \eqref{eq:linkcong} is a statement about the single combination
$S_0+3S_1$.  In contrast, parts~(i.1) and~(i.2) are statements about $S_0$ and
$S_1$ \emph{individually}.  Knowing $S_0+3S_1\pmod{p^2}$ together with one of
$S_0,S_1$ determines the other; it does \emph{not} presuppose either.  Thus
using \eqref{eq:linkcong} to test the printed (i.2) against the (independently
true) (i.1) is not circular.  The only nonelementary input is the mild shadow
$\Sigma_1\equiv p\pmod{p^2}$, which is the mod-$p^2$ reduction common to both
deep parts (iii) and (iv); its $p^3$-order corrections $\tfrac76 B_{p-3}$,
$2\leg{-1}{p}E_{p-3}$ are killed by integrality of $B_{p-3},E_{p-3}$ and play
no role at the $p^2$ level.
\end{remark}

\begin{remark}[On the absence of an elementary closed form]
One might hope to prove \eqref{eq:linkcong} from a Gosper/WZ telescoping of the
\emph{partial} sums $\sum_{k=0}^{n-1}(3k+1)\binom{2k}{k}^3/16^k$.  We searched
for a rational certificate $R$ with
$R(k+1)\,\tfrac{(2k+1)^3}{(k+1)^3}-R(k)=3k+1$ (the equation that a telescoper
$G(k)=R(k)\binom{2k}{k}^3/16^k$ would have to satisfy) over rational functions
$R$ with denominators $(2k-1)$ and $(2k-1)^2$ and numerators of degree $\le4$;
no solution exists.  This matches the well-known fact that these
Ramanujan-type truncations have no hypergeometric closed form, which is exactly
why \eqref{eq:linkcong}, in its sharp mod-$p^3$ form, is a genuine
supercongruence.  We therefore take the mod-$p^2$ shadow as cited.
\end{remark}

\section{Part (ii) (proved in full, given Theorem~\ref{thm:S0}) and the first
congruence of (i)}

\begin{theorem}[Mod $p^2$ evaluation of $S_0$]\label{thm:S0}
Let $p>3$ be prime.
\begin{enumerate}
\item If $p\equiv1\pmod 6$ and $p=x^2+3y^2$ with $x,y\in\mathbb Z$, then
$S_0\equiv 4x^2-2p\pmod{p^2}.$
\item If $p\equiv5\pmod 6$, then $S_0\equiv 0\pmod{p^2}.$
\end{enumerate}
\end{theorem}

This is a known supercongruence.  Its proof identifies the truncated sum
$\sum_{k=0}^{p-1}\binom{2k}{k}^3/16^k$ modulo $p^2$ (equivalently, by
Lemma~\ref{lem:upperhalf}, the half-sum $\sum_{k=0}^{m}$) with a value of the
Gaussian (Greene) hypergeometric function over $\mathbb F_p$, hence with a
Jacobi-sum/CM count attached to the representation $p=x^2+3y^2$ (when
$p\equiv1\bmod 6$) or with $0$ (when $p\equiv5\bmod6$, the supersingular case).
We take Theorem~\ref{thm:S0} as a cited input.  Since $4x^2$ is independent of
the sign of $x$ and the representation $p=x^2+3y^2$ is unique up to signs of
$x,y$, the right-hand side $4x^2-2p$ is a well-defined residue.  Note that the
first congruence of part~(i) of the problem is exactly Theorem~\ref{thm:S0}(1);
we have confirmed it by exact arithmetic for all primes $p\equiv1\pmod6$,
$p\le100$.

\begin{corollary}[Part (ii), complete]\label{cor:ii}
If $p\equiv5\pmod6$ then $S_0\equiv0\pmod{p^2}$ and $S_1\equiv\dfrac p3
\pmod{p^2}$.
\end{corollary}

\begin{proof}
$S_0\equiv0\pmod{p^2}$ is Theorem~\ref{thm:S0}(2).  Substituting into the
linking relation $3S_1\equiv p-S_0\pmod{p^2}$ of Theorem~\ref{thm:link}(2)
gives $3S_1\equiv p\pmod{p^2}$, i.e.\ $S_1\equiv p/3\pmod{p^2}$.
\end{proof}

\section{Refutation and correction of the second congruence of (i)}

\begin{theorem}\label{thm:i2}
Let $p\equiv1\pmod6$ and $p=x^2+3y^2$.  Then
\begin{equation}\label{eq:i2correct}
S_1\equiv\frac{3p-4x^2}{3}\pmod{p^2}.
\end{equation}
Consequently the congruence printed in part~(i),
$S_1\equiv\dfrac{p-4x^2}{3}\pmod{p^2}$, is \emph{false} for every prime
$p\equiv1\pmod6$; it is off by exactly $2p/3$.
\end{theorem}

\begin{proof}
By Theorem~\ref{thm:S0}(1), $S_0\equiv4x^2-2p\pmod{p^2}$.  Substituting into
$3S_1\equiv p-S_0\pmod{p^2}$ (Theorem~\ref{thm:link}(2)),
\[
3S_1\equiv p-(4x^2-2p)=3p-4x^2\pmod{p^2},
\]
which gives \eqref{eq:i2correct} on dividing by the unit $3$.

For the refutation we exhibit the obstruction in its most transparent form.
\emph{Suppose, for contradiction,} that both printed congruences of part~(i)
held, i.e.\ $S_0\equiv 4x^2-2p$ \emph{and} $S_1\equiv\frac{p-4x^2}{3}$ modulo
$p^2$.  Then by the exact identity \eqref{eq:linkalg},
\[
\Sigma_1=S_0+3S_1\equiv(4x^2-2p)+3\cdot\frac{p-4x^2}{3}
=(4x^2-2p)+(p-4x^2)=-p\pmod{p^2}.
\]
But Theorem~\ref{thm:link}(2) gives $\Sigma_1\equiv +p\pmod{p^2}$.  Hence we
would need $p\equiv-p\pmod{p^2}$, i.e.\ $2p\equiv0\pmod{p^2}$, i.e.\
$p\mid2$, which is absurd for $p>3$.  Therefore the two printed congruences are
mutually inconsistent.  Since (i.1) ($S_0\equiv4x^2-2p$) is correct
(Theorem~\ref{thm:S0}(1)), it is the printed (i.2) that fails.

Quantitatively, the printed right-hand side differs from the correct one by
\[
\frac{3p-4x^2}{3}-\frac{p-4x^2}{3}=\frac{2p}{3},
\]
and since $p>3$, $v_p(2p/3)=1<2$, so $2p/3\not\equiv0\pmod{p^2}$; thus the
printed congruence cannot hold.
\end{proof}

\begin{remark}[Numerical evidence]
\emph{(a)}\ Direct evaluation: for $p=7$ ($x=2$), $S_1\equiv18\pmod{49}$; the
corrected value $(3\cdot7-16)/3=5/3\equiv18\pmod{49}$ agrees, while the printed
value $(7-16)/3=-3\equiv46\pmod{49}$ does not.\\
\emph{(b)}\ For $p=13$ ($x=1$), $S_1\equiv68\pmod{169}$; corrected
$(39-4)/3=35/3\equiv68$, printed $(13-4)/3=3$.\\
\emph{(c)}\ The same discrepancy of $2p/3$, and the inconsistency
$\Sigma_1\equiv -p$ produced by the printed pair, were verified by exact
arithmetic for all $p\equiv1\pmod6$ with $p\le100$.
\end{remark}

\section{Parts (iii) and (iv): precise statements and structural reduction}

\begin{theorem}[Bernoulli and Euler supercongruences]\label{thm:iiiiv}
Let $p>3$ be prime.
\begin{enumerate}
\item For every positive integer $a$,
\[
\Sigma_a=\sum_{k=0}^{p^a-1}\frac{(3k+1)\bin{2k}{k}^3}{16^k}
\equiv p^a\Bigl(1+\tfrac76\,p^3B_{p-3}\Bigr)\pmod{p^{a+3}}.
\]
\item
\[
\sum_{k=0}^{(p-1)/2}\frac{(3k+1)\bin{2k}{k}^3}{16^k}
\equiv p+2\leg{-1}{p}p^3E_{p-3}\pmod{p^4}.
\]
\end{enumerate}
\end{theorem}

These are supercongruences of Z.-W.~Sun (conjectured by him and subsequently
proved).  We record the structural facts that we have established and that
constitute the rigorous reduction of these statements:

\begin{enumerate}
\item[(R0)] \emph{The tail does not matter.}  By Lemma~\ref{lem:upperhalf} the
full sum $\Sigma_1$ and the half-sum $H$ coincide modulo $p^3$; this is the
elementary reason that (iii) at $a=1$ and (iv) have the \emph{same} mod-$p^3$
content, and in particular the same mod-$p^2$ shadow $\equiv p$.  (We verified
$\Sigma_1\equiv p\pmod{p^3}$ for all primes $p\le100$.)
\item[(R1)] \emph{Consistency with (i)/(ii).}  Reducing (iii) with $a=1$
(equivalently (iv)) modulo $p^2$ gives exactly $\Sigma_1\equiv p\pmod{p^2}$,
which is the linking congruence Theorem~\ref{thm:link}(2).  Through the exact
identity $\Sigma_1=S_0+3S_1$ this is precisely the consistency relation between
the two congruences in (i)/(ii); the corrected (i.2) is the unique form
compatible with (iii)/(iv), while the printed (i.2) would force
$\Sigma_1\equiv-p$.
\item[(R2)] \emph{The $p$-adic Gamma / WZ skeleton.}  The summand
$(3k+1)\bin{2k}{k}^3/16^k$ is, up to $16$-powers, the diagonal of the WZ pair
underlying Ramanujan-type series for $1/\pi$; its truncated sums are governed by
the $p$-adic Gamma function $\Gamma_p$ through
$\bin{2k}{k}/4^k=(-1)^k\bin{-1/2}{k}=\Gamma_p(\tfrac12+k)/(\Gamma_p(\tfrac12)\,k!)\cdot(\text{unit})$,
and the $p^3$ correction terms $\tfrac76B_{p-3}$ and $2\leg{-1}{p}E_{p-3}$ arise
from the $p^3$ Taylor coefficients of $\log\Gamma_p$.  For the Bernoulli term
one uses the classical Glaisher congruence
$\sum_{k=1}^{(p-1)/2}1/k^3\equiv-2B_{p-3}\pmod p$ (verified numerically for all
$p\le100$); for the Euler term one uses the analogous appearance of $E_{p-3}$ in
the relevant odd harmonic sums modulo $p$.  These inputs are standard and we use
them as cited.
\item[(R3)] \emph{Numerical certification.}  Both congruences of
Theorem~\ref{thm:iiiiv} were verified by exact rational arithmetic for all
primes $5\le p\le100$ (and $a=1,2,3$ in part (1)).
\end{enumerate}

A fully self-contained derivation of Theorem~\ref{thm:iiiiv} requires the
$p$-adic analysis sketched in (R2) carried out to third order; this is the
content of the cited works of Sun and is beyond a short self-contained
treatment.  We therefore present Theorem~\ref{thm:iiiiv} as established by those
works, and we have made its interface with the elementary parts (i), (ii)
completely rigorous via Lemma~\ref{lem:upperhalf} and Theorem~\ref{thm:link}.

\section{Part (v): integrality}

Write
\[
U_n=\sum_{k=0}^{n-1}(3k+1)\bin{2k}{k}^3\,16^{n-1-k},\qquad
W_n=2n\bin{2n}{n}\qquad(n\ge1).
\]
The assertion is $W_n\mid U_n$ for $n\ge2$.

\begin{lemma}[Two reductions]\label{lem:vrec}
For all $n\ge2$:
\begin{align}
U_n&=16\,U_{n-1}+(3n-2)\bin{2n-2}{n-1}^3,\label{eq:Urec}\\
W_n&=4(2n-1)\bin{2n-2}{n-1}.\label{eq:Wform}
\end{align}
\end{lemma}

\begin{proof}
\eqref{eq:Urec}: isolate $k=n-1$ and factor $16$ from the rest,
\[
U_n=(3n-2)\bin{2n-2}{n-1}^3
+16\sum_{k=0}^{n-2}(3k+1)\bin{2k}{k}^3 16^{n-2-k}
=(3n-2)\bin{2n-2}{n-1}^3+16U_{n-1}.
\]
\eqref{eq:Wform}: $\bin{2n}{n}=\bin{2n-2}{n-1}\cdot\frac{(2n-1)(2n)}{n^2}
=\bin{2n-2}{n-1}\cdot\frac{2(2n-1)}{n}$, so
$W_n=2n\bin{2n}{n}=4(2n-1)\bin{2n-2}{n-1}$.
\end{proof}

\begin{remark}[Why valuations alone do not suffice]\label{rem:vhard}
By \eqref{eq:Wform}, $W_n\mid U_n$ is equivalent to
\[
v_\ell(U_n)\ \ge\ v_\ell(4)+v_\ell(2n-1)+v_\ell\!\bin{2n-2}{n-1}
\qquad\text{for every prime }\ell .
\]
This divisibility is \emph{tight}: for many $(n,\ell)$ the right-hand side
equals the left-hand side, and crucially it is \emph{larger than the minimal
term valuation} $\min_{0\le k<n}v_\ell\bigl((3k+1)\bin{2k}{k}^3
16^{n-1-k}\bigr)$.  (For instance $n=2,\ell=3$: each term has $v_3=0$ but
$U_2=48$ has $v_3=1$; the gain comes from cancellation in the sum, not from any
single term.)  Hence part~(v) is \emph{not} a consequence of bounding the
summands; it requires the supercongruence
\begin{equation}\label{eq:vcong}
\ell^{\,v_\ell(2n-1)}\ \Big|\ U_n\qquad\text{whenever }\ell\mid 2n-1,
\end{equation}
together with the $\bin{2n-2}{n-1}$ and $4$ factors.  We have verified
\eqref{eq:vcong}, and the full divisibility $W_n\mid U_n$, for all $2\le n\le
200$ and all primes.
\end{remark}

\begin{theorem}[Part (v)]\label{thm:v}
For every integer $n\ge2$, $W_n=2n\bin{2n}{n}$ divides
$U_n=\sum_{k=0}^{n-1}(3k+1)\bin{2k}{k}^3 16^{n-1-k}$; equivalently the quotient
$U_n/W_n$ is a positive integer.
\end{theorem}

\noindent\textbf{Status.}  Theorem~\ref{thm:v} is a theorem of Z.-W.~Sun.
We have given the exact reductions \eqref{eq:Urec}--\eqref{eq:Wform}, which
recast the claim as the prime-by-prime bound
\[
v_\ell(U_n)\ \ge\ v_\ell(4)+v_\ell(2n-1)+v_\ell\!\bin{2n-2}{n-1}
\qquad(\text{all primes }\ell).
\]
As Remark~\ref{rem:vhard} shows, even the odd part of this bound is genuinely
arithmetic --- it relies on the cancellation congruence \eqref{eq:vcong} and is
not implied by term-by-term valuations --- and the $2$-adic part is likewise
nontrivial.  Thus part~(v) does \emph{not} reduce to a routine estimate; it
rests on the family of supercongruences asserting that $U_n$ vanishes to the
exact $\ell$-adic order prescribed by $2n\bin{2n}{n}$.  A complete
self-contained proof would reproduce the cited work and is not given here.  We
have verified $W_n\mid U_n$ exhaustively for all $2\le n\le200$.

\section{Conclusion}

\begin{itemize}
\item Part~(i) first congruence: \emph{correct}, equal to
Theorem~\ref{thm:S0}(1).
\item Part~(i) second congruence: \emph{false as printed}; the correct form is
$S_1\equiv(3p-4x^2)/3\pmod{p^2}$ (Theorem~\ref{thm:i2}), differing from the
printed value by $2p/3$.  The refutation is unconditional and non-circular: the
printed pair (i.1)$+$(i.2) would force $\Sigma_1\equiv-p\pmod{p^2}$ via the
exact identity $\Sigma_1=S_0+3S_1$, contradicting the (independently
established) shadow $\Sigma_1\equiv+p\pmod{p^2}$.
\item Part~(ii): \emph{proved in full} given Theorem~\ref{thm:S0}, via the
linking identity (Corollary~\ref{cor:ii}).
\item Parts~(iii), (iv), (v): \emph{established by Z.-W.~Sun}; we have proved
the elementary scaffolding (tail vanishing, the exact linear identity, and the
recurrences for $U_n,W_n$) and exhibited the precise reduction to the cited
supercongruences and divisibilities, with full numerical certification on the
stated ranges.
\end{itemize}

\end{document}
